% GENERATED FILE. Edit canonical lessons, then run npm run generate:books.
% canonical-source-hash: fnv1a64:696a3845164185be
% canonical-lessons: KA-C516-vrtti, KA-C516-hudde, KA-C516-arhate, KA-C516-tajna, KA-C516-parinati

\chapter{Careers}
\label{ch:ka-careers}
\hlchaptermodality{516}

\begin{chapteropening}
\textbf{By the end of this chapter:} \emph{I can say a profession, a post, a qualification, an expert and expertise.}

Complete the last lesson of chapter 516: I can say a profession, a post, a qualification, an expert and expertise.
\end{chapteropening}

\section[\texorpdfstring{\kn{ವೃತ್ತಿ} (vṛtti)}{ವೃತ್ತಿ (vṛtti)}]{\kn{ವೃತ್ತಿ} (vṛtti) — a profession; a career}

\label{lesson:KA-C516-vrtti}



\begin{quote}
Before the new one: say the Kannada for to groan, then the Kannada for to droop.
\end{quote}

\subsection*{You'll want to know: \kn{ವೃತ್ತಿ}}
\textbf{\kn{ವೃತ್ತಿ}} — \emph{vṛtti} — \textquotedblleft{}a profession; a career\textquotedblright{}.

\begin{quote}
\kn{ಅವನು} \kn{ಶಿಕ್ಷಕ} \kn{ವೃತ್ತಿಯನ್ನು} \kn{ಆರಿಸಿದನು}.
\end{quote}

\begin{quote}
\emph{avanu śikṣaka vṛttiyannu ārisidanu.}
\end{quote}

\begin{quote}
He chose the teaching profession.
\end{quote}

\kn{ವೃತ್ತಿಪರ}, professional.

\begin{cousinweb}
\textbf{\kn{ವೃತ್ತಿ}} (\emph{vṛtti}) is borrowed from Sanskrit \emph{vṛtti}, \textquotedblleft{}rolling, course; mode of life, livelihood\textquotedblright{} (Monier-Williams), from the root \emph{vṛt}, \textquotedblleft{}to turn\textquotedblright{}.
\end{cousinweb}

\begin{grammarlens}[title={What you've built}]
One more word for describing what you lived and felt.
\end{grammarlens}

\subsection*{Your turn}
\begin{itemize}
\raggedright
  \item \emph{Say it:} \emph{vṛtti}
  \item \emph{Say it:} \emph{vṛtti}, once more
  \item \emph{Say it:} \emph{avanu śikṣaka vṛttiyannu ārisidanu.} — He chose the teaching profession.
  \item \emph{Recall:} say the Kannada for to groan, then the Kannada for to droop, then say \emph{vṛtti} again
\end{itemize}

\begin{tcolorbox}[breakable,colback=teal!4,colframe=teal!35!black,title={Before you move on}]
What does \kn{ವೃತ್ತಿ} mean? (\textquotedblleft{}A profession; a career\textquotedblright{}.) Say it once more.
\end{tcolorbox}
\section[\texorpdfstring{\kn{ಹುದ್ದೆ} (hudde)}{ಹುದ್ದೆ (hudde)}]{\kn{ಹುದ್ದೆ} (hudde) — a post; a position}

\label{lesson:KA-C516-hudde}



\begin{quote}
Before the new one: say the Kannada for to droop, then the Kannada for a profession.
\end{quote}

\subsection*{You'll want to know: \kn{ಹುದ್ದೆ}}
\textbf{\kn{ಹುದ್ದೆ}} — \emph{hudde} — \textquotedblleft{}a post; a position\textquotedblright{}.

\begin{quote}
\kn{ಅವಳಿಗೆ} \kn{ಬ್ಯಾಂಕಿನಲ್ಲಿ} \kn{ಒಳ್ಳೆಯ} \kn{ಹುದ್ದೆ} \kn{ಸಿಕ್ಕಿತು}.
\end{quote}

\begin{quote}
\emph{avaḷige byāṅkinalli oḷḷeya hudde sikkitu.}
\end{quote}

\begin{quote}
She got a good post at the bank.
\end{quote}

A rank or seat in an office.

\begin{cousinweb}
\textbf{\kn{ಹುದ್ದೆ}} (\emph{hudde}) is usually traced to Hindustani \emph{ohdā}, \textquotedblleft{}an office, a rank\textquotedblright{}, itself from Arabic; it came in with the language of offices.
\end{cousinweb}

\begin{grammarlens}[title={What you've built}]
One more word for describing what you lived and felt.
\end{grammarlens}

\subsection*{Your turn}
\begin{itemize}
\raggedright
  \item \emph{Say it:} \emph{hudde}
  \item \emph{Say it:} \emph{hudde}, once more
  \item \emph{Say it:} \emph{avaḷige byāṅkinalli oḷḷeya hudde sikkitu.} — She got a good post at the bank.
  \item \emph{Recall:} say the Kannada for to droop, then the Kannada for a profession, then say \emph{hudde} again
\end{itemize}

\begin{tcolorbox}[breakable,colback=teal!4,colframe=teal!35!black,title={Before you move on}]
What does \kn{ಹುದ್ದೆ} mean? (\textquotedblleft{}A post; a position\textquotedblright{}.) Say it once more.
\end{tcolorbox}
\section[\texorpdfstring{\kn{ಅರ್ಹತೆ} (arhate)}{ಅರ್ಹತೆ (arhate)}]{\kn{ಅರ್ಹತೆ} (arhate) — a qualification; eligibility}

\label{lesson:KA-C516-arhate}



\begin{quote}
Before the new one: say the Kannada for a profession, then the Kannada for a post.
\end{quote}

\subsection*{You'll want to know: \kn{ಅರ್ಹತೆ}}
\textbf{\kn{ಅರ್ಹತೆ}} — \emph{arhate} — \textquotedblleft{}a qualification; eligibility\textquotedblright{}.

\begin{quote}
\kn{ಈ} \kn{ಕೆಲಸಕ್ಕೆ} \kn{ಬೇಕಾದ} \kn{ಅರ್ಹತೆ} \kn{ಅವನಿಗೆ} \kn{ಇದೆ}.
\end{quote}

\begin{quote}
\emph{ī kelasakke bēkāda arhate avanige ide.}
\end{quote}

\begin{quote}
He has the qualification this job needs.
\end{quote}

\kn{ಅರ್ಹ}, worthy, eligible + -\kn{ತೆ}, -ness.

\begin{cousinweb}
\textbf{\kn{ಅರ್ಹತೆ}} (\emph{arhate}) is built from Sanskrit \emph{arha}, \textquotedblleft{}meriting, deserving, worthy\textquotedblright{} (Monier-Williams), with the noun ending -\emph{tā}, Kannada -\kn{ತೆ}.
\end{cousinweb}

\begin{grammarlens}[title={What you've built}]
One more word for describing what you lived and felt.
\end{grammarlens}

\subsection*{Your turn}
\begin{itemize}
\raggedright
  \item \emph{Say it:} \emph{arhate}
  \item \emph{Say it:} \emph{arhate}, once more
  \item \emph{Say it:} \emph{ī kelasakke bēkāda arhate avanige ide.} — He has the qualification this job needs.
  \item \emph{Recall:} say the Kannada for a profession, then the Kannada for a post, then say \emph{arhate} again
\end{itemize}

\begin{tcolorbox}[breakable,colback=teal!4,colframe=teal!35!black,title={Before you move on}]
What does \kn{ಅರ್ಹತೆ} mean? (\textquotedblleft{}A qualification; eligibility\textquotedblright{}.) Say it once more.
\end{tcolorbox}
\section[\texorpdfstring{\kn{ತಜ್ಞ} (tajña)}{ತಜ್ಞ (tajña)}]{\kn{ತಜ್ಞ} (tajña) — an expert; a specialist}

\label{lesson:KA-C516-tajna}



\begin{quote}
Before the new one: say the Kannada for a post, then the Kannada for a qualification.
\end{quote}

\subsection*{You'll want to know: \kn{ತಜ್ಞ}}
\textbf{\kn{ತಜ್ಞ}} — \emph{tajña} — \textquotedblleft{}an expert; a specialist\textquotedblright{}.

\begin{quote}
\kn{ಅಪ್ಪನನ್ನು} \kn{ಹೃದಯ} \kn{ತಜ್ಞರ} \kn{ಬಳಿ} \kn{ಕರೆದುಕೊಂಡು} \kn{ಹೋದೆವು}.
\end{quote}

\begin{quote}
\emph{appanannu hṛdaya tajñara baḷi karedukoṇḍu hōdevu.}
\end{quote}

\begin{quote}
We took father to a heart specialist.
\end{quote}

Plural: \kn{ತಜ್ಞರು}, \emph{tajñaru}.

Often after a field: \kn{ಹೃದಯ} \kn{ತಜ್ಞ}, a heart specialist.

\begin{cousinweb}
\textbf{\kn{ತಜ್ಞ}} (\emph{tajña}) is borrowed from Sanskrit \emph{tajjña}, \textquotedblleft{}knowing that; a man who knows\textquotedblright{} (Monier-Williams).
\end{cousinweb}

\begin{grammarlens}[title={What you've built}]
One more word for describing what you lived and felt.
\end{grammarlens}

\subsection*{Your turn}
\begin{itemize}
\raggedright
  \item \emph{Say it:} \emph{tajña}
  \item \emph{Say it:} \emph{tajña}, once more
  \item \emph{Say it:} \emph{appanannu hṛdaya tajñara baḷi karedukoṇḍu hōdevu.} — We took father to a heart specialist.
  \item \emph{Recall:} say the Kannada for a post, then the Kannada for a qualification, then say \emph{tajña} again
\end{itemize}

\begin{tcolorbox}[breakable,colback=teal!4,colframe=teal!35!black,title={Before you move on}]
What does \kn{ತಜ್ಞ} mean? (\textquotedblleft{}An expert; a specialist\textquotedblright{}.) Say it once more.
\end{tcolorbox}
\section[\texorpdfstring{\kn{ಪರಿಣತಿ} (pariṇati)}{ಪರಿಣತಿ (pariṇati)}]{\kn{ಪರಿಣತಿ} (pariṇati) — expertise; mastery}

\label{lesson:KA-C516-parinati}



\begin{quote}
Before the new one: say the Kannada for a qualification, then the Kannada for an expert.
\end{quote}

\subsection*{You'll want to know: \kn{ಪರಿಣತಿ}}
\textbf{\kn{ಪರಿಣತಿ}} — \emph{pariṇati} — \textquotedblleft{}expertise; mastery\textquotedblright{}.

\begin{quote}
\kn{ಅವರಿಗೆ} \kn{ಕಂಪ್ಯೂಟರ್} \kn{ಕೆಲಸದಲ್ಲಿ} \kn{ಪರಿಣತಿ} \kn{ಇದೆ}.
\end{quote}

\begin{quote}
\emph{avarige kampyūṭar kelasadalli pariṇati ide.}
\end{quote}

\begin{quote}
She has expertise in computer work.
\end{quote}

\kn{ಪರಿಣತ}, an expert.

\begin{cousinweb}
\textbf{\kn{ಪರಿಣತಿ}} (\emph{pariṇati}) is borrowed from Sanskrit \emph{pariṇati}, from \emph{pari-ṇam}, \textquotedblleft{}to change, ripen\textquotedblright{} (Monier-Williams): skill that has ripened.
\end{cousinweb}

\begin{grammarlens}[title={What you've built}]
One more word for describing what you lived and felt.
\end{grammarlens}

\subsection*{Your turn}
\begin{itemize}
\raggedright
  \item \emph{Say it:} \emph{pariṇati}
  \item \emph{Say it:} \emph{pariṇati}, once more
  \item \emph{Say it:} \emph{avarige kampyūṭar kelasadalli pariṇati ide.} — She has expertise in computer work.
  \item \emph{Recall:} say the Kannada for a qualification, then the Kannada for an expert, then say \emph{pariṇati} again
\end{itemize}

\begin{tcolorbox}[breakable,colback=teal!4,colframe=teal!35!black,title={Before you move on}]
What does \kn{ಪರಿಣತಿ} mean? (\textquotedblleft{}Expertise; mastery\textquotedblright{}.) Say it once more.
\end{tcolorbox}
